\begin{abstract}
eBPF lets applications execute custom logic in the kernel for tasks such as packet processing and system observability. However, the existing compilation pipeline fails to select efficient implementations available on the target CPU for some computations. Exploiting these opportunities requires existing approaches to trade off native optimization capability against the complexity of trusted code. Bytecode rewriting is constrained by the existing JIT's instruction-selection rules, whereas adding native optimizations requires trusting the corresponding program transformations and code generation. We present \tool, an extensible eBPF compilation framework that enables more efficient native implementations while reusing the existing verifier for safety analysis. The framework identifies optimization opportunities in userspace and selects native implementations provided by the kernel through extended instructions. Each extension describes its behavior using eBPF instructions for safety analysis within the complete program and executes using a semantically equivalent native implementation. We implement seven categories of local optimizations on x86-64 and ARM64. Relative to the stock Linux eBPF JIT, microbenchmarks achieve geometric mean speedups of 1.24$\times$ and 1.22$\times$, respectively. Throughput increases by 7.4\% for Cilium on x86-64 and by 7.3\% for Katran on ARM64.
\end{abstract}
